\subsubsection{Determinación constructiva y función del registro dodecafásico}
\label{alpha-ampl-dependencias}

La definición anterior permite precisar qué significa determinar alfa desde
la estructura. El dominio de la operación contiene cuatro registros de
procedencia conocida: las tres secuencias regionales y el registro integral
de fase. La normalización incorpora, además, una condición de frontera
compatible con la prolongación. Una vez fijados esos datos, la división
euclídea determina cada resto y su cociente. La proposición
\ref{prop:alpha-normalizacion} extiende esa determinación al sistema
completo. La independencia respecto de un valor físico objetivo es, por
tanto, una propiedad del orden de construcción: el valor que se compara al
final no interviene en la selección de ninguno de los argumentos de esta
aplicación.

Esta independencia causal debe distinguirse de la independencia algebraica
entre números. La primera se establece identificando el dominio, las
operaciones y la procedencia de los coeficientes; la segunda sería una
afirmación sobre relaciones polinómicas. El resultado utilizado aquí es el
primero. La construcción determina una coordenada mediante registros
producidos por APP--TRIT--TPK y permite comprobar posteriormente las
relaciones que satisface. Ese orden hace posible utilizar una coordenada ya
generada en una operación ulterior sin convertirla en un parámetro externo.

El registro $K$ tiene una función más precisa que la de una corrección
numérica añadida a una fórmula. Sus doce posiciones proceden de la
transformación reversible del registro firmado. La ecuación
$U=\mathcal H_{12}K$ conserva la posibilidad de volver a ese registro,
mientras las diferencias transversales y la suma total permiten otra
reconstrucción. En la evaluación periódica, el orden de las posiciones
determina el peso posicional de cada componente. El mismo conjunto de
componentes, dispuesto en otro orden, produce en general otra coordenada
racional. Por ello el origen de fase y la orientación forman parte de la
definición del número que interviene en \eqref{eq:alpha-identidad-completa}.

La normalización de alfa compone esa información con las regiones. Su
identidad local distingue dos operaciones sucesivas. Primero se forma el
balance orientado de bloques; después se representa ese balance mediante
un resto canónico y un cociente que enlaza posiciones contiguas. La primera
operación conserva el signo del balance y permite la extracción posterior
de un soporte. La segunda permite formar los prefijos de la coordenada
real. Ambas lecturas permanecen relacionadas porque el estado conserva los
datos anteriores y posteriores a la normalización. El soporte firmado y la
expansión posicional son, así, realizaciones diferentes de una composición
aritmética explícita.

La identidad telescópica expresa el contenido global de esta composición.
Los cocientes interiores se cancelan por pares al evaluar un prefijo,
mientras sus extremos permanecen en la igualdad. Esa cancelación es una
propiedad de la suma ponderada; no elimina los cocientes del estado ni
permite suprimirlos antes de efectuar la suma. La frontera derecha de una
ventana recibe la contribución de la continuación. Por esa razón, la
evaluación de una ventana aislada y la restricción de una sección
prolongada son operaciones distintas. El cambio del bloque final $663$ a
$662$, demostrado anteriormente, ofrece una realización explícita de esa
diferencia sin alterar los bloques ya estabilizados a su izquierda.

En este sentido, la bidireccionalidad tiene una formulación aritmética
concreta. La generación proporciona prolongaciones compatibles; la
normalización transmite la información entera de su frontera desde las
posiciones de menor peso hacia las de mayor peso. La ecuación conserva la
contribución de ambas direcciones mediante sus índices y sus condiciones
de contorno. La reversibilidad en la fibra fijada se comprueba reconstruyendo
$K$ a partir de los demás registros y de los cocientes completos. Esa
inversión certifica la consistencia de la operación original; la dirección
genealógica continúa siendo la que produce primero el registro y después
la coordenada de alfa.

La segunda determinación tiene un contenido complementario. En ella, el
objeto inmediato es un balance de vacancias y el parámetro aparece como
su raíz en un dominio aislante. La vía posicional proporciona una
normalización canónica de registros; la vía analítica estudia la anulación
de un operador. La exposición de ambas requiere conservar precisamente
esa diferencia de funciones. Sus relaciones con el estado común han de
expresarse mediante las aplicaciones de refinamiento y la comparación de
sus realizaciones completas, tal como se desarrolla en la subsección
correspondiente. El acuerdo de una aproximación finita constituye un
control localizado; la identificación de las construcciones completas
tiene el dominio y el alcance de su propio enunciado.

Finalmente, esta organización fija la posición de alfa en el artículo. Su
ecuación compone la generación regional de la sección~\ref{sec:generacion}
con el registro \(K\) de la sección~\ref{sec:k-registro}. La aplicación
electrónica del Artículo I incorpora después \(\alpha_{\mathrm{HMT}}\) como
factor de su sección central; el presente artículo recibe de esa aplicación
las secciones de acción \eqref{v:ant:eq:el-secciones-accion} y las emplea en
la normalización gravitatoria. La prolongación excepcional de la
sección~\ref{exc:seccion} recibe, en cambio, el registro, sus soportes y la
incidencia conservada. Ambas continuaciones comparten procedencia y realizan
operaciones diferentes: evaluación de una sección central y construcción de
una configuración de frontera. Sus mapas internos mantienen visible la
estructura completa que determina la coordenada escalar.
